\originalchapter{来源对应与范围边界}
23份教师 PDF 共149个实际页面; 第24讲为复习, 官网没有讲义. 各文件页数、原始直链、SHA256及许可检查见 source-manifest.json, 这些原件不在中文源码包中. 对应关系如下; “择要”明确表示未覆盖整份讲义, 不能据此宣称逐页全译.
\begin{longtable}{@{}p{.12\textwidth}p{.32\textwidth}p{.49\textwidth}@{}}
\toprule 原讲次&原主题&中文册与边界\\\midrule\endhead
1--4&群、子群、陪集、循环群&第1--2章, 完整核心论证, 合并重复例子\\
5--6&置换与共轭&第3章, 完整核心论证\\
7--10&同构、同态、商及同构定理&第2章, 完整核心论证\\
11&交错群&第3章符号, 作业6中补完整单性证明\\
12&群表示与小阶群&第3章二面体表示, 模拟期末第3题分类至10阶; 一般自由群构造未展开\\
13&Sylow&第4章完整证明; 作业和模拟卷完整公开题面题解\\
14--16&环、基本性质、同态与理想&第5章, 完整核心论证\\
17--20&分式域、素/极大理想、特殊及欧几里得整环&第5--6章, 完整核心论证\\
21&多项式&第7章及作业11, 完整核心论证\\
22&快速素性检验&第14章择要: Fermat、多项式判别、平方根见证; 不含 AKS 参数界与完整复杂度证明, 不声称完整覆盖\\
23&群作用与自同构&群作用在第4章, 域自同构第11--13章, 多项式自同构在作业11\\
\bottomrule
\end{longtable}
RES.18-011第1--6讲的基础群论并入第1--3章; 第17--19、21--23讲的作用、稳定子、共轭与 Sylow 并入第3--4章. 其第7--11讲线性代数仅作为模论应用的先修参考, 不逐讲重写. 第12--16、20、24--35讲的正交几何、离散群、二十面体、双线性/厄米空间、谱定理、线性Lie群及 Hilbert 第三问题未纳入.

RES.18-012第8--13讲的环、理想、分解并入第5--7章; 第19--22讲的模、关系矩阵、Smith 与分解构成第8--9章; 第24--27、29--31讲的域、扩张、有限域、Galois 构成第10--13章. 第23讲 Noether 环仅使用 PID 升链停止的自足特例, 未证明一般 Hilbert 基定理. 第1--7讲及附录的群表示与特征标、第14--18讲数域理想分解与类群、第28讲函数域几何、第32--35讲一般根式求解与判别式理论未完整纳入. 本原元、Artin 定理及各核心定理中的省略步骤由本册补写; 不是英文整页替换为中文标题.
原18.703第22讲关于模奇合数的平方根1数量的引理漏掉素数幂情形, 第14章以模9反例和正确条件修正. 原作业8把环直积当作一般环余积的附注不成立, 第5章给反例. 作业2第3题需排除平凡群并澄清“没有非平凡真子群”; 作业6第12题的交错群单性陈述需 $n\ge3$. 作业11在文件标题中误写作业9, 本册按官网作业顺序编号11并保留映射. 这些改动是数学勘正, 不宣称官方已勘误.

同源合订PDF与单讲PDF内容重复, 学生原 TeX 包与合订PDF是同一笔记的不同格式, 故不另算学科或重复抄入正文. 英文官方包只作来源归档, 原图不复用; 本册唯一域格图由 TikZ 独立绘制.
